\documentclass{newlayout}
%Bitte hier den enstprechenden Ort einsetzen z.B. Braunschweig und die Akademienummer
\Akademie{Braunschweig}{2015}{1}

\usepackage[english, ngerman]{babel}
\usepackage{misc}
\usepackage{multicol}

%\usepackage{amsmath}%wird automatisch durch newlayout.cls geladen
\usepackage{amssymb}
\usepackage{amsfonts}
\usepackage{tikz}
\usepackage{dsfont}
\usepackage{blindtext}
\usepackage{tipa}
\usepackage{algpseudocode}
\usepackage{url}
\def\UrlBreaks{\do\a\do\b\do\c\do\d\do\e\do\f\do\g\do\h\do\i\do\j\do\k\do\l%
\do\m\do\n\do\o\do\p\do\q\do\r\do\s\do\t\do\u\do\v\do\w\do\x\do\y\do\z\do\0%
\do\1\do\2\do\3\do\4\do\5\do\6\do\7\do\8\do\9\do\-\do\_\do\/\do\%}
\urlstyle{same}

% hinzugefügt, um Fehler 'pdfTeX error (font expansion): auto expansion is only possible with scalable' zu vermeiden
\usepackage{lmodern}
\setkomafont{descriptionlabel}{\normalfont\bfseries}
\addtokomafont{paragraph}{\normalfont}
\usepackage{footnote}
\usepackage[flushmargin,hang,ragged]{footmisc}
\deffootnote{1em}{1em}{%
\textsuperscript{\thefootnotemark\ }
}
%\setlength{\abovedisplayskip}{5pt}
%\setlength{\belowdisplayskip}{5pt}

\newcommand{\media}{media/}
\newcommand{\norm}[1]{\left\lVert#1\right\rVert}

%%%%%Mathe-Definitionen
\newtheorem{Def}{Definition}
\newtheorem{Sat}{Satz}
\newtheorem{Bew}{Beweis}
\newtheorem{Thm}{Theorem}

\setlength\abovedisplayshortskip{0pt}
\setlength\belowdisplayshortskip{0pt}
\setlength\abovedisplayskip{3pt}
\setlength\belowdisplayskip{3pt}
%%%%Ende Mathe-Definitionen

\begin{document}
%\input{titel}
\setcounter{page}{3}

\setcounter{tocdepth}{1}
\tableofcontents

\setcounter{secnumdepth}{1}

\setcounter{page}{7}
\setcounter{chapter}{0}

%Angabe, bis zu welcher Stufe die sections im Text nummeriert werden sollen.
\settocdepth{2}

\course{2}{Learning from Data}%%% 
\begin{coursetitle}
  \centerline{Learning from Data} 
  \bigskip
  \Large \centerline{Mathematische Theorien und Methoden für das digitale Zeitalter}
  \bigskip
 \includegraphics[width=\columnwidth]{\media LearningFromData_Logo}
 \label{fig:titelbild}
  \bigskip
\end{coursetitle}
%\begin{dsafigure}
%\begin{center}
%\includegraphics[width=.9\columnwidth]{Titelbild-fehlt.png}
%\caption{meine Bildunterschrift}
%label{fig:meinbild}
%\end{center}
%\end{dsafigure}

%Einleitung
\section{Aus der Kursbeschreibung}
\authors{Philipp Moritz KL, Fanny Yang KL}

Test.

Jeder Kurs beginnt mit einem Teil "`1 Aus der Kursbeschreibung"', entnommen aus dem Programmheft.
Hier sollte nur der den Inhalt des Kurses beschreibende
Teil der Kursbeschreibung erscheinen, nicht
die Erwartungen an die Teilnehmenden.




Gegen Ende des Kurses können die Teilnehmenden je nach Interesse weitere Themen diskutieren. Es kann zum Beispiel das ${H_2}^+$-Ion als ein
\section{Einleitung}
\authors{Lukas Kamm}
\begin{quote}
\glqq Okay, Google, wie komme ich von hier am schnellsten zum Alexanderplatz? \grqq\\
\\
\glqq Geh' zu Fuß 100 Meter zur U-Bahn-Haltestelle Moritzplatz und nimm die U8 in Richtung Osloer Straße. Die nächste U-Bahn fährt in 6 Minuten.\grqq
\end{quote}
Ein Szenario wie es vermutlich jeder Smartphone-Benutzer kennt. Die Spracheingabe wird verwendet, um seinem Gerät Fragen zu stellen: Nach dem Wetter, aktuellen Aktienkursen oder dem Weg von A nach B. Doch wie ist es möglich, Sprache zu erkennen? Sie setzt sich immerhin komplex aus verschiedenen Komponenten wie den Lauten, der Tonhöhe, Betonung und Unterschieden in der Lautstärke zusammen, verläuft zeitlich und variiert dazu von Sprecher zu Sprecher. Die Antwort lautet: \emph{maschinelles Lernen}. Das Prinzip von Lernalgorithmen ist es, den Computer anhand verschiedener \emph{Trainingsbeispiele} darauf zu trainieren, wie etwa das Wort \glqq Alexanderplatz\grqq klingt, um es im Umkehrschluss im Gesprochenen wiederzuerkennen. Diese Lernalgorithmen haben sich in den vergangenen zehn Jahren in der Spracherkennung etabliert, spielen aber auch in anderen Feldern eine bedeutende Rolle.\\
Ein anderes Beispiel für die Benutzer traditioneller Kommunikationsmedien: Die Deutsche Post stellt täglich rund 64 Millionen Briefe zu (\url{https://www.deutschepost.de/de/q/qualitaet_gelb.html}). Da bei dieser Menge an zu identifizierenden Postleitzahlen, Straßen und Orten der Mensch zu langsam ist, kommen dort Lernalgorithmen in der Texterkennung von gedruckter und handgeschriebener Schrift zum Einsatz, um die Anschrift des Empfängers digital zu erfassen und die Sendungen ins richtige Fach zu transportieren.\\
Ein verbreitetes Prinzip der künstlichen Intelligenz ist das künstliche neuronale Netz. Der Mensch bedient sich hierbei an der Natur, indem die Informationsverarbeitung im Nervensystem des menschlichen Gehirns und Rückenmarks mathematisch beschrieben und auf Computern simuliert wird.\\
Erste Ansätze zu neuronalen Netzen gab es in den 1940er-Jahren. Bald zeigten sich jedoch in der Praxis Grenzen dieser Konzepte mangels effizienter Algorithmen und der geringen Rechenleistung damaliger Rechenanlagen. Erst in den 1970er- und 1980er-Jahren, beispielsweise durch die Entwicklung der sogenannten \emph{Backpropagation} (Fehlerrückführung) und der wachsenden Leistungsfähigkeit von Rechnern, konnten neuronale Netze praktikabel in der Erkennung von Text, Sprache, Gesichtern, Mustern in der medizinischen Diagnostik, etwa bei einer fMRT-Aufzeichnung oder in der Zeitreihenanalyse von Wetter oder Aktien, eingesetzt werden.\\
Im Kurs möchten wir zunächst ein Verständnis für die mathematische Modellierung verschiedener Lernalgorithmen gewinnen. Dazu eignen wir uns anfangs Grundlagen der höheren Mathematik, konkret aus den Bereichen lineare Algebra, Analysis, Wahrscheinlichkeitstheorie und der Optimierung, an. Anschließend lernen wir \emph{Support Vector Machines} kennen. Ein Algorithmus, der darauf trainiert wird, ein Datum anhand verschiedener Eigenschaften in eine richtige Klasse einzuordnen. Des Weiteren betrachten wir künstliche neuronale Netze und die Umsetzung und Anwendung beider in Software.\\
Zunächst beginnen wir mit den mathematischen Grundlagen.

%\documentclass{article}

%\usepackage[ngerman]{babel}
%\usepackage[utf8]{inputenc}
%\usepackage{color}
%\usepackage[a4paper,lmargin={2cm},rmargin={2cm},
%tmargin={2.5cm},bmargin = {2.5cm}]{geometry}
%\usepackage{amssymb}
%\usepackage{amsmath}
%\usepackage{graphicx}
%\usepackage{multicol}
%\usepackage{amsthm}

%\setlength{\columnsep}{1cm}
%\newtheorem{defi}{Definition}
%\newtheorem*{defi*}{Definition}

%\begin{document}
%\begin{multicols*}{2}

\section{Vektoren}
\authors{Vivien Thi, Luca Bohn}
In der linearen Algebra nehmen Vektoren eine zentrale Rolle ein. Ein Vektor wird in der Schule als eine Menge an Pfeilen gelehrt, die parallel, gleichgerichtet und gleich lang sind. Betrachten wir jedoch nicht nur den dreidimensionalen Raum $\mathbb{R}^3$, sondern $\mathbb{R}^n$ f\"ur beliebige $n$, so ist die Vorstellung eines Pfeils nicht immer m\"oglich. Dies bedingt die Notwendigkeit einer anderen Definition.

Ein \textit{Vektor} ist eine geordnete Aufz\"ahlung von Objekten und wird auch \textit{Tupel} $(a_1, a_2,\dots,a_n)$ genannt. Entscheidend bei Tupeln ist die Reihenfolge der Objekte, die auch mehrfach vorkommen k\"onnen. Dies steht im Kontrast zum Begriff einer Menge. Die Addition von Vektoren $u, v\in\mathbb{R}^n$ ist dann definiert durch
\begin{equation*}
u + v = \begin{pmatrix} u_1 \\ u_2 \\ \dots \\ u_n \end{pmatrix} + \begin{pmatrix} v_1 \\ v_2 \\ \dots \\ v_n \end{pmatrix} = \begin{pmatrix} u_1 + v_1 \\ u_1+v_2 \\ \dots \\ u_1+v_n \end{pmatrix}
\end{equation*} 
und die skalare Multiplikation eines Vektors $u\in\mathbb{R}^n$ mit $\lambda\in\mathbb{R}$ ist definiert durch
\begin{equation*}
\lambda u = \lambda \begin{pmatrix} u_1 \\ u_2 \\ \dots \\ u_n \end{pmatrix} = \begin{pmatrix} \lambda u_1 \\ \lambda u_2 \\ \dots \\ \lambda u_n \end{pmatrix}.
\end{equation*}



\subsection{Vektor- \& Unterraum}
Vektoren bilden die Elemente eines \textit{Vektorraumes V}. Addieren wir zwei Vektoren eines Vektorraumes oder multiplizieren wir sie mit einem Skalar, so muss die Summe bzw. das Produkt ebenfalls ein Element des Vektorraumes sein. Das hei\ss t
\begin{enumerate}
\item $u, v \in V \text{ }\Rightarrow u+v \in V$
\item $u\in V, \lambda \in \mathbb{R} \text{ }\Rightarrow \lambda u \in V$
\end{enumerate}
\vspace{10pt}
Ein \textit{Unterraum} $U$ ist eine Teilmenge eines Vektorraumes $V$. Es gelten die oben genannten Eigenschaften eines Vektorraumes.
\begin{enumerate}
\item $\text{ }u, v \in U,\text{ } u+v \in U$
\item $\text{ }u \in U, \lambda \in \mathbb{R} ,\text{ }\lambda u \in U$
\end{enumerate}



\subsection{Span}
Der span$(v_1, v_2,\dots,v_k)$ ist die Menge aller Vektoren, die mit den Linearkombinationen der Vektoren $v_1, v_2, \dots, v_k$ darstellbar sind, das hei\ss t
\begin{align*}
\text{span}(v_1,\dots,v_k) = \{u \in V : u = \lambda_1v_1+\dots+\lambda_kv_k\} \\
\text{ mit } \lambda_1,\dots,\lambda_k \in \mathbb{R}
\end{align*}



\begin{theorem}
Ein Span ist immer auch ein Vektorraum.
\end{theorem}
\begin{proof} 
Angenommen, $v \in \text{span}(v_1,\dots,v_k)$,
dann wissen wir durch die Definition des Spans, dass man $u$ und $v$ auch darstellen kann als
\begin{align*}
u = \sum\limits_{i=1}^{k} \lambda_iv_i, \quad v = \sum\limits_{i=1}^{k} \alpha_iv_i
\end{align*}
Wegen der Addition bei Vektorr\"aumen gilt dann
\begin{align*}
	u + v = \sum\limits_{i=1}^{k} \lambda_iv_i + \sum\limits_{i=1}^{k} \alpha_iv_i= \sum\limits_{i=1}^{k} (\alpha_i+\lambda_i)v_i
\end{align*}
Wegen der Multiplikation bei Vektorr\"aumen gilt dann
\begin{align*}
	\lambda\cdot v = (\lambda \cdot \lambda_1) v_1+\dots+(\lambda \cdot \lambda_k)v_k
\end{align*}
Dadurch hat man gezeigt, dass ein Span die Eigenschaften eines Vektorraumes erf\"ullt.
\end{proof}



\subsection{Basis \& Dimension}
Jeder Vektorraum kann durch eine Menge linear unabh\"angiger Basisvektoren dargestellt werden. Diese Menge wird \textit{Basis} eines Vektorraumes genannt. Es gibt beliebig viele unterschiedliche Basen. Die Anzahl der Basisvektoren bestimmt die \textit{Dimension p} eines Vektorraums.



\subsection{Lineare Unabh\"angigkeit}
Die Vektoren $v_1,\dots,v_k$ mit $v_i \in V$ sind \textit{linear unabh\"angig}, falls $\lambda_iv_i+\dots+\lambda_kv_k = 0$ nur f\"ur $\lambda_i=\dots=\lambda_k= 0$ gilt. 

%\end{multicols*}
%\end{document}

%\documentclass{article}

%\usepackage{tikz}
%\usepackage[utf8]{inputenc}
%\usepackage[ngerman]{babel}
%\usepackage{color}
%\usepackage[a4paper,lmargin={2cm},rmargin={2cm},
%tmargin={2.5cm},bmargin = {2.5cm}]{geometry}
%\usepackage{amssymb}
%\usepackage{amsmath}
%\usepackage{graphicx}
%\usepackage{multicol}
%\usepackage{amsthm}

%\newtheorem*{defi*}{Definition}
%\newtheorem*{defi2*}{Definition}

%\begin{document}
%\begin{multicols*}{2}

% \newcommand{\norm}[1]{\left\lVert#1\right\rVert}

\subsection{Norm und Skalarprodukt}

\begin{Def}
$\norm{v}$ ist eine Norm im Vektorraum $V$ falls
\begin{enumerate}
\item 
\begin{enumerate}
\item $\norm{v} \geq 0$ für alle $v \in V$
\item $\norm{v} = 0$ nur für $v = 0$
\end{enumerate}
\item $\norm{u+v} \leq \norm{u} + \norm{v}$ für alle $v \in V$ (Dreiecks-Ungleichung)
\item $\norm{\lambda v} = |\lambda| \norm{v}$ für alle $\lambda \in \mathbb{R}$ und $v \in V$ (Homogenit\"at)
\end{enumerate}
\end{Def}

In Worte gefasst ergeben sich die folgenden Eigenschaften: Eine Norm kann nie negativ sein, ist die Norm 0 bedeutet das, dass der Vektor der Nullvektor ist. Wenn man zwei Vektoren
addiert, ist die Norm des resultierenden Vektors kleiner oder gleich der Summe der Normen der beiden Ausgangsvektoren. Außerdem ist die
Norm eines Vektors, der mit einem reellen Faktor multipliziert wurde, gleich dem Produkt aus der Norm des Vektors und dem 
Betrag des Faktors.

\begin{Def}
$f:V\times V \rightarrow\mathbb{R}$ ist ein Skalarprodukt, wenn\\
\begin{enumerate}
\item $\left\langle  \alpha u + \beta v, w\right\rangle = \alpha \left\langle  u,w\right\rangle + \beta \left\langle  v,w\right\rangle $ (Bilinearität)
\item $\left\langle  u,v\right\rangle = \left\langle  v,u\right\rangle $ für alle $u,v \in V$
\item $\left\langle  v,v\right\rangle \geq 0$ mit $\left\langle  v,v\right\rangle = 0$ nur für $v = 0$
\end{enumerate}
\end{Def}
Wir können ein Skalarprodukt im euklidischen Raum $\mathbb{R}^n$ mit
\begin{equation}
\left\langle u,v\right\rangle  = u_1 v_1 + \dots + u_n v_n
\label{skal}
\end{equation}
definieren. An der ausgeschriebenen Form lässt sich überprüfen, ob es sich dabei wirklich um ein Skalarprodukt handelt. Hierbei kann  $\left\langle \alpha u + \beta v, w\right\rangle $  umgeformt werden. Durch das Ausmultiplizieren der Produkte und Sortieren der Summe in Terme mit $\alpha$ oder  $\beta$, gelangen wir zu der aus der Definition geforderten Form 
\begin{align*}
\langle \alpha u + \beta v, w\rangle &= \sum_{i=1}^n (\alpha u_i + \beta v_i)w_i \\
&= \alpha \left\langle u,w\right\rangle  + \beta \left\langle v,w\right\rangle .
\end{align*}
Die zweite Bedingung können wir mit dem Kommutativgesetz beweisen, wenn wir dieses auf die ausgeschriebene Form anwenden.
Für die dritte Bedingung bekommen wir $\sum_{i=1}^n v_i v_i = \sum_{i=1}^n v_i^2$. Da durch das Quadrieren keiner der Summanden kleiner als 0 werden kann und auch nur für $v_i = 0$ genau 0 werden kann, ist auch diese Bedingung erfüllt.

\begin{dsafigure}
\centering
\begin{tikzpicture}
\coordinate (A) at (0,0);
\coordinate (B) at (4,1.5);
\coordinate (C) at (1,3);
\draw[thick, ->] (A) -- (B) node[pos=0.5,sloped,below]{$v$};
\draw[thick, ->] (A) -- (C) node[pos=0.5,sloped, above]{$u$};
\draw[dashed, ->] (B) -- (C) node[pos=0.5,sloped, above]{$v-u$};
\draw (1,0.375) arc (21:75:1);
\node[] at (45:0.7) {$\theta$};
\end{tikzpicture}
\caption {Anwendung des Kosinussatzes}
\label{Kosinussatz}
\end{dsafigure}

Da grundsätzlich $\langle v,v\rangle  = \norm{v}^2$ gilt, kann über die Abbildung \ref{Kosinussatz} auch die Aussage getroffen werden, dass $\norm{v-u}^2 = \left\langle v-u,v-u\right\rangle $. Mit der Bilinearität des Skalarproduktes und der eben getroffenen Aussage können wir die Gleichung zu 
$\lVert v-u\rVert ^2 = \lVert v\rVert ^2 - 2\left\langle v,u\right\rangle  + \lVert u\rVert ^2$
umformen.
Zusätzlich folgt aus dem Kosinussatz, dass $\norm{v-u}^2 = \norm{u}^2 + \norm{v}^2 - 2 \norm{u} \norm{v} \cos\theta$. Setzt man beide Terme gleich, gelangt man über einige wenige Umformungen zu der Form $\left\langle u,v\right\rangle  = \norm{u} \norm{v} \cos\theta$, was ebenfalls eine Möglichkeit ist, das Skalarprodukt darzustellen.
Da sich der Betrag des Kosinus nur zwischen 0 und 1 bewegt, folgt aus dieser Gleichung direkt die Cauchy-Schwarz Ungleichung
\begin{equation}
\left\langle u,v\right\rangle  \leq \norm{u} \norm{v}.
\label{cauchy}
\end{equation}

Eine häufig genutzte Norm ist die Euklidische Norm, welche durch $\norm{v}_2 = \sqrt{v_1^2 + \dots + v_n^2}$ definiert ist und auch durch die obrige Definition des Skalarproduktes \eqref{skal} induziert wird.

Die erste und die dritte Bedingung aus der Definition für eine Norm zeigen sich aus der ausgeschriebenen Form der Euklidischen Norm. Da alle Komponenten dabei quadriert werden und die Wurzel gezogen wird, muss die erste Bedingung stimmen. Ein Faktor $\lambda$, der ebenfalls quadriert in der Wurzel steht kann aus der Summe ausgeklammert werden und als Faktor vor die Wurzel gelangen. Der Betrag ergibt sich dabei daraus, dass $\lambda$ nach dem Vorgang nur positiv sein kann, auch wenn es zuvor negativ war.

Um zu beweisen, dass die Dreiecksungleichung $\norm{u+v} \leq\norm{u} + \norm{v}$  gilt, benötigen wir die Dreiecksungleichung für reelle Zahlen und die Cauchy-Schwarz Ungleichung. Wir gehen dabei von der quadrierten Form des linken Terms aus und erhalten
\begin{align*}
\norm{u+v}^2 &= |\left\langle u+v,u+v\right\rangle|\\
&\leq |\left\langle u,u\right\rangle| + 2|\left\langle u,v\right\rangle| + |\left\langle v,v\right\rangle|\\
&\leq \norm{v}^2 + 2\norm{u} \norm{v} + \norm{v}^2\\
&= (\norm{u} + \norm{v})^2.
\end{align*}

Man sieht sofort, dass der Ausgangsterm kleiner oder gleich dem Ergebnis ist, womit bewiesen ist, dass die Euklidische Norm eine Norm ist.

%\end{multicols*}
%\end{document}

\subsection{Lineare Abbildungen}
In diesem Abschnitt werden wir die elementarste Form einer Funktion kennenlernen, die einen Vektorraum auf einen anderen Vektorraum abbildet. Diese Funktionen nennt man \textit{lineare Abbildungen}.
\begin{Def}
	\label{Def:Def_1}
Eine Funktion $f:U \stackrel{}{\rightarrow} V$ ist linear, wenn folgende Bedingungen gelten: \\ 
1. $f(u+v) = f(u)+f(v) \ \forall \ u,v \in U$\\
2. $f(\lambda v) = \lambda f(v) \ \forall \ v \in U \ \text{und} \ \lambda \in \mathbb{R}$
\end{Def}

Das besondere an solchen Abbildungen ist, dass sie von \textit{Matrizen} eindeutig definiert werden. Wie Matrizen definiert sind und wie wir sie verwenden können, lernen wir im nächsten Kapitel. Im Folgenden gehen wir davon aus, dass $U$ die Dimension $n$ hat und $V$ die Dimension $m$. 

\subsubsection{Matrix}
Eine Matrix ist eine spezielle Anordnung von $n\cdot m$ Zahlen in $n$ Zeilen und $m$ Spalten. Man kann eine Matrix aber auch als eine Zusammenfassung von $m$ Vektoren mit $n$ Komponenten ansehen, wobei ein Vektor eine Spalte bildet. Die Komponente in der $i$-ten Zeile und $j$-ten Spalte wird als $a_{ij}$ geschrieben und die ganze Spalte mit $a_j$ notiert.\\
\begin{equation*}
A= \left(
   \begin{array}{cccc}
	  a_{11} & a_{12} & ... & a_{1m}\\
		a_{21} & a_{22} & ... & a_{2m}\\
		...   & ...   & ... & ...  \\
		a_{n1} & a_{n2} & ... & a_{nm}
	 \end{array}
	 \right) = (a_1, a_2, ..., a_m)\\
\end{equation*}
\\
Da wir nun wissen was eine Matrix ist, führen wir in dem folgenden Abschnitt zwei neue Operationen ein: die \textit{Matrix-Vektor-Multiplikation} und \textit{Matrix-Matrix-Multiplikation}.

\subsubsection{Matrix-Vektor-Multiplikation}
Bei einer Matrix-Vektor-Multiplikation wird eine Matrix $A \in  \mathbb{R}^{n \times m}$ mit einem Vektor $v \in \mathbb{R}^m$ multipliziert. Dabei bildet jede Multiplikation zwischen einer Matrixzeile und einem Vektor ein Skalarprodukt. Eine Multiplikation zwischen Matrizen und Vektoren, bei denen die Anzahl nicht übereinstimmt, ist nicht definiert. Deshalb ist es wichtig, dass die Matrix so viele Spalten wie der Vektor Dimensionen hat.
\begin{align*}
	Av &= 
	\left(
   \begin{array}{cccc}
	  a_{11} & a_{12} & ... & a_{1m}\\
		a_{21} & a_{22} & ... & a_{2m}\\
		...   & ...   & ... & ...  \\
		a_{n1} & a_{n2} & ... & a_{nm}
	 \end{array}
	\right) 
	\cdot 
	\left(
	 \begin{array}{c}
	  v_1\\
		v_2\\
		...\\
		v_m
	 \end{array}
	\right)\\
	&= \left(
	 \begin{array}{c}
	 a_{11} \cdot v_1 + a_{12} \cdot v_2 + ... + a_{1m} \cdot v_m\\
	 a_{21} \cdot v_1 + a_{22} \cdot v_2 + ... + a_{2m} \cdot v_m\\
		...\\
	 a_{n1} \cdot v_1 + a_{n2} \cdot v_2 + ... + a_{nm} \cdot v_m\\
	 \end{array}
	\right)
\end{align*}

Allgemein lässt sich die Matrix-Vektor-Multiplikation folgendermaßen berechnen

\begin{equation*}
(Av)_i = \sum_{j=1}^{n}a_{ij}v_j ,
\end{equation*}

wobei $(AV)_i$ der $i$-te Eintrag des Vektors $n$ ist, der das Ergebnis der Multiplikation ist.\\
Eine Eigenschaft der linearen Abbildung ist, dass sich jede Abbildung als ein solches Matrix-Vektor-Produkt schreiben lässt. 

\subsubsection{Matrix-Matrix-Multiplikation}
Eine Matrix-Matrix-Multiplikation ist nichts anderes als mehrere Matrix-Vektor-Multiplikationen hintereinander. Haben wir zwei Matrizen $A \in \mathbb{R}^{n \times m}$ und $B \in \mathbb{R}^{m \times p}$ und berechnen $AB$, nehmen wir jede Spalte von $B$ als Vektor und multiplizieren ihn mit $A$ wie oben beschrieben. Die Vektoren, die man als Ergebnisse erhält, fasst man wieder in einer Matrix mit der Dimension $\mathbb{R}^{n \times p}$ zusammen. Allgemein können wir die Matrix-Matrix-Multiplikation folgendermaßen berechnen
\begin{equation*}
	(AB)_{il} = \sum_{n=i}^{m}a_{ij} \cdot b_{jl}
\end{equation*}

\subsubsection{Komposition}
\begin{Thm}
Eine Komposition $h(x) = (f \circ g)(x)$ aus den linearen Abbildungen $f$ und $g$, ist auch eine lineare Abbildung.
\end{Thm}

\begin{proof}
Da $f$ und $g$ linear sind gelten die Bedingungen aus Definition \ref{Def:Def_1}:
\begin{align*}
	h(u+v) &= (f \circ g)(u+v)\\
	&= f(g(u+v))\\
	&= f(g(u)+g(v))\\
	&= f(g(u)) + f(g(v))\\
	&= (f\circ g)(u) + (f\circ g)(v) = h(u) + h(v)
\end{align*}
Des Weiteren gilt:
\begin{align*}
	h(\lambda v) &= (f \circ g)(\lambda v)\\
	&=f(g(\lambda v))\\
	&= f(\lambda g(v))\\
	&=\lambda f(g(v))\\
	&=\lambda (f \circ g)(v) = \lambda h(v)
\end{align*}
Da beide Bedingungen aus Definition \ref{Def:Def_1} für $h$ erfüllt sind, ist auch eine Komposition aus zwei linearen Abbildungen linear.
\end{proof}




\section{Analysis}
  

Analysis ist ein wichtiges Teilgebiet der Mathematik. Sie ist Grundlage für Optimierung und somit auch von großer Bedeutung für das Maschinelle Lernen. Im Kurs war für uns besonders die Stetigkeit und Differenzierbarkeit von ein- und mehrdimensionalen Funktionen wichtig.

\subsection{Stetigkeit und Differenzierbarkeit}

Wir beginnen mit einigen theoretischen Grundlagen der Analysis.
\begin{Def} Die Zahl $L$ heißt \emph{Grenzwert} der Funktion $f(x)$ für $x$ gegen $a$, wenn es für alle $\epsilon > 0$ ein $\delta > 0$ existiert, sodass für alle $|x-a| < \delta$ gilt $|f(x)-L|<\epsilon$
\end{Def}
\begin{Def}
Eine Funktion $f: D \subseteq \mathbb{R}\rightarrow \mathbb{R}$ ist \emph{stetig}, wenn $\lim_{w\rightarrow x}f(w)=f(x)$ auf ihrem gesamten Definitionsbereich gilt, das heißt sie ist stetig in $x$ für alle $x\in D$.
\end{Def}
\begin{Def}
Eine Funktion $f:\mathbb{R}\rightarrow \mathbb{R}$ ist an einer Stelle $x$ \emph{differenzierbar}, wenn eine lineare Abbildung $ l_ x (h)$ und ein Restglied $r(h)$ existieren, sodass $f(x+h)=f(x)+l_{x}(h)+r(h)$ wobei $\lim_{h\rightarrow 0} r(h)/h= 0$ gilt.
\end{Def}
Wir betrachten eine einfache Folgerung aus diesen Definitionen.

\begin{exercise}Zeige, dass eine Funktion $f:\mathbb{R} \rightarrow \mathbb{R}$, welche differenzierbar an der Stelle $x \in \mathbb{R}$ ist dort auch stetig ist.
\end{exercise}
Lösung: Wir überlegen zuerst, welche Bedingungen gegeben sind. Eine Funktion ist an der Stelle $x$ differenzierbar, wenn $f(x+h)=f(x)+l_x(h)+ r(h)$  
mit $\lim_{h \rightarrow 0} r(h)/h= 0$. Um die Stetigkeit nachzuweisen müssen wir zeigen, dass $\lim_{w \rightarrow x} f(w)=f(x)$. Mit $w = x+h$ ist dies äquivalent zu $\lim_{h \rightarrow 0} f(x)+l_x(h)+r(h)=f(x)$. Es genügt also zu zeigen, dass  $l_x(h)$ und $r(h)$ gegen Null gehen.
\begin{itemize}
  \item Würde $r(h)$ nicht gegen Null gehen, so würde $r(h)/h$ nicht gegen Null gehen. Widerspruch zur Differenzierbarkeit.
  \item Der lineare Term $l_x(h)$ kann auch als $f'(x) \cdot h$ dargestellt werden. Da $h$ gegen Null geht folgt $l_x(h)\rightarrow 0$.
\end{itemize}
Insgesamt folgt $\lim_{w\rightarrow x} f(w) = f(x)$.
\subsection{Gradient}
\begin{Def}
Der \emph{Gradient} $\nabla f(x)$ einer Funktion $f:\mathbb{R}^n\rightarrow \mathbb{R}$ ist ein Vektor, wobei jeder Eintrag die partielle Ableitung der entsprechenden Komponente ist, also
\begin{equation*} \begin{split} \nabla f(x) = \left( \begin{array}{c}
\partial_{x_1} f(x) \\
\partial_{x_2} f(x) \\
... \\
\partial_{x_n} f(x)
\end{array}
\right).
\end{split} \end{equation*} 
\end{Def}
Der Gradient zeigt immer in Richtung des steilsten Anstiegs einer Funktion. Der negative Gradient in die Richtung des steilsten Abstiegs, weshalb wir das Gradientenabstiegsverfahren verwenden können, um das Minimum einer mehrdimensionalen Funktion zu finden. Wichtig für das Gradientenabstiegsverfahren ist die folgende Aufgabe.
\begin{exercise}
Zeige, dass es immer möglich ist für eine differenzierbare Funktion $f$ einen Skalar $\lambda$ zu finden, sodass
\begin{equation}\label{lambda-ungleichung}f(x + \lambda \nabla f(x)) \leq f(x)\end{equation} 
\end{exercise}
Lösung:
Schreibe die Ungleichung \eqref{lambda-ungleichung} um in
\begin{equation*}
f(x) + \text{D} f(x)h + r(h) \leq f(x)
\end{equation*}
Hierbei ist $h = \lambda \nabla f(x)$ und $r(h)$ das Restglied.
%% Das Ganze kann man noch weiter umschreiben, indem man das Skalarprodukt bildet: \\ \\
%% $f(x) + \lambda \langle \nabla f(x), \nabla f(x) \rangle + r(h) \leq f(x)$ \\ \\
%% Anschließend kann man von der Regel Gebrauch machen, dass sich ein Skalarprodukt, das in der Form $\langle u,u \rangle$ vorliegt, in die Ausprägung der Euklidischen Norm von $||u||^2$ umgewandelt werden kann. \\
%% Daraus folgt: \\ \\
%% $f(x) + \lambda \cdot || \nabla f(x)||^2 + r(h) \leq f(x)$ \\ \\
%% Um nun die obige Ungleichung (1) zu beweisen, muss man zeigen, dass \\ $\lambda ||\nabla f(x)||^2 + r(h) \leq 0$ ist.\\
%% Dafür wird $\lambda$ ausgeklammert. Anschließend erhält man:

%% \begin{equation} \lambda (||\nabla f(x)||^2 + \frac{r(h)}{\lambda} \end{equation}

%% Da $||\nabla f(x)||$ quadriert wird, wird dieser Term automatisch positiv. Da man den gesamten Term aber negativ haben will, wählt man $\lambda < 0.$ Schließlich kennt jeder die Regel, dass ein Produkt negativ wird, falls eine ungerade Zahl von Faktoren negativ ist. Da wir $\lambda$ negativ gewählt haben, muss \\ $||\nabla f(x)||^2$$(2) \lambda (||\nabla f(x)||^2 + \frac{r(h)}{\lambda}$ positiv sein.\\
%% Bei $||\nabla f(x)||^2||$ stellt dies aus den eben genannten Gr\"unden kein Problem dar. Deswegen muss nur noch gezeigt werden, dass es mindestens einen Wert von $r(h)$ gibt, der gr\"oßer oder gleich 0 ist. Ansonsten würde der Wert in der Klammer negativ werden k\"onnen und damit die unspr\"ungliche Aussage (1) widerlegen.\\
%% Damit wir zeigen k\"onnen, dass $r(h) \geq 0$ ist, formten wir zuerst den Term (2) um:\\

%% $\lambda (||\nabla f(x)||^2 + \frac{r(h)}{\lambda})$ \\

%% $\Leftrightarrow \lambda(||\nabla f(x)||^2 + \frac{r(\lambda \nabla f(x))}{\lambda})$ \\

%% Bei dieser Umformung wurde h wieder mit $\lambda\nabla f(x)$ ersetzt.

%% $\leftrightarrow \lambda ||\nabla f(x)|| \cdot (||\nabla f(x)|| + \frac{r(\lambda \nabla f(x))}{\lambda || f(x)||})$


%% Durch diese Umformungen schlussfolgerten wir, dass f\"ur die Funktion $h(x)$ folgendes gelten muss:\\
%% \begin{equation} h(\lambda)=||\nabla f(x)|| + \frac{r(\lambda\nabla f(x))}{\lambda||\nabla f(x)||} \geq 0 \end{equation}


%% Da in $||\nabla f(x)$ kein $\lambda$ steht ist dieser Term konstant. Gleichzeitig wird durch das Bilden der euklidischen Norm festgelegt, dass ein Wert größe/gleich null vorliegt. \\ %der Satz ist noch nicht wiklich schön formuliert
%% Deswegen muss nur noch gezeigt werden, dass dder zweite Summand ebenfalls positiv werden kann. Dazu bildeten wir den Grenzwert. \\ %Die Zeitformen stimmen sehr oft nicht
%% $\lim\limits_{h \rightarrow 0}{\frac{r(\lambda \nabla f(x))}{\lambda ||\nabla f(x)||}} = \frac{r(\lambda \nabla f(x))}{||\lambda  \nabla f(x)||}$ \\ %hier kommt noch kein h vor. Umschreiben?

%% Man kann $\lambda \nabla f(x)$ wieder durch h ersetzen, woraus folgt: \\
%% $\lim\limits_{h \rightarrow 0} \frac{r(h)}{h}$ \\

%% Der Grenzwert des Restbetrages strebt gegen Null. Daraus lässt sich schließen, dass es einen Punkt gibt, ab welchem die Ungleichung der Funktion $h(\lambda)$ (3) erfüllt ist und damit auch die anfängliche Aussage (1). \\ %ist hier h(\lambda) und (3) beies nötig?

\subsection{Summen-, Produkt- und Kettenregel}

Falls mehrere Funktonen miteinander verknüpft sind gelten folgende drei Regeln.
\begin{itemize}
\item Die \emph{Summenregel} besagt, dass der Gradient der Funktion $h(x)=f(x)+g(x)$ die Summe der Gradienten ist, das heißt
\begin{equation*} D(f + g)(x)=\text{D} f(x)+\text{D} g(x)\end{equation*}

\item Gegeben seien die differenzierbaren Funktionen $f:\mathbb{R}^k\rightarrow\mathbb{R}^n$, $g:\mathbb{R}^n\rightarrow\mathbb{R}^k$ und $h(x)=f(g(x))$. Die \emph{Kettenregel} besagt, dass
\begin{equation*}
\text{D}h(x) = \text{D}f(g(x)) \cdot \text{D} g(x).
\end{equation*}

\item Gegeben zwei differenzierbaren Funktionen $f:\mathbb{R}^m\rightarrow\mathbb{R}$ und $g:\mathbb{R}^m\rightarrow\mathbb{R}$, so gilt für das Produkt der Funktionen $h(x)=f(x) \cdot g(x)$ die \emph{Produktregel} \begin{equation*}\text{D}h(x)=f(x) \cdot \text{D}g(x)+\text{D}f(x) \cdot g(x).\end{equation*}
 \end{itemize}
Wir beweisen nun die Produktregel.
Da die Funktionen $f$ und $g$ differenzierbar sind gilt
\begin{align*}
  f(x+h) &= f(x) +  Df(x) h + r_{f}(h)\\
  g(x+h) &= f(x) + Dg(x) h + r_{g}(h)
\end{align*}
wobei $\lim_{h\rightarrow 0} r_f(h)/h = 0$ und $\lim_{h\rightarrow 0} r_g(h)/h = 0$.
Wir schreiben zuerst die Ableitung von $h(x+v)$ als Kombination aus den Ableitungen von $f(x)$ und $g(x)$ auf, also $h(x+v) =  f(x+v)\cdot g(x+v)$.
Das Produkt wird ausmultipliziert und wir erhalten
\begin{equation*}
\begin{split} h(x+v) & = f(x) \cdot g(x) \\ & + f(x)  \cdot D_{g}(x) \cdot v \\ & + f(x) \cdot r_{g}(v) \\ & + Df(x) \cdot v \cdot g(x) \\ & +  Df(x) \cdot v \cdot D_{g}(x) \cdot v  \\ & + Df(x) \cdot v \cdot r_{g}(v) \\ & +  r_{g}(v) \cdot g(x) \\ & + r_{g}(x) \cdot Dg(x) \cdot v \\ & + r_{g}(v) \cdot r_{g}(v) \end{split} \end{equation*}   %Die Formel passt nicht in eine Zeile. Umbruch?

Von den ganzen Summanden müsssen letztendlich alle eliminiert werden, bis auf 
\begin{equation*} f(x) \cdot g(x) + f(x) \cdot Dg(x) \cdot v + Df(x) \cdot v \cdot g(x) \end{equation*} 
Dafür müssen wir zeigen, dass die übrigen Summanden gegen Null gehen. Um das zu erreichen bilden wir den Grenzwert
\begin{align*}
\lim \limits_{v \rightarrow 0} &||v||^{-1}(||f(x) \cdot r_{g}v + Df(x) \cdot v \cdot Dg(x) \cdot v +{}\\
&\qquad + Df(x) \cdot v \cdot r_{g}(v) r_{f}(v) \cdot g(x) +{}\\
&\qquad + r_{g}(v) \cdot Dg(x) \cdot v + r_{g}(v) \cdot r_{f}(v)||)
\end{align*}
Damit die Summanden einzeln behandelt werden können, wenden wir die Dreiecksungleichung an. Dann können wir von jedem Summanden einzeln den Limes bilden, um zu zeigen, dass der gesamte Grenzwert gegen Null geht. Dies ist aber nur der Fall, falls der Limes jedes Summanden gegen Null geht. Wir zeigen das exemplarisch für 

\begin{equation*}
\lim\limits_{v \rightarrow 0} \frac{||Df(x) \cdot r_{g}(v)||}{||v||}
\end{equation*}

$Df(x)$ ist eine Matrix mit einer Zeile und m Spalten in der jeweils ein Eintrag $\leq c$ steht. Daraus folgern wir:
\begin{align*}
&\lim\limits_{v \rightarrow 0} \frac{m \cdot c \cdot ||v|| \cdot ||r_{g}(v)||}{||v||}\\
&\qquad= m \cdot c \cdot \lim\limits_{v \rightarrow 0} \frac{ ||v|| \cdot ||r_{g}(v)||}{||v||} = 0
\end{align*} 
Damit haben wir gezeigt, dass der Summand gegen Null geht. Ähnlich behandelt man auch die anderen Summanden und auch diese gehen alle gegen Null.
Die Produktregel ist damit bewiesen. 

















\section{Wahrscheinlichkeitstheorie}


\vspace{15pt}


1933 ver\"offentlichte der russische Mathematiker Andrey Kolmogorov sein Buch \textit{Foundations of the Theory of Probability}, in dem er die drei Axiome der Wahrscheinlichkeitstheorie aufstellte. Es existiert ein Wahrscheinlichkeitsraum, der sich aus den drei Elementen $\Omega, \mathcal{F}, \mathds{P}$ zusammensetzt.

\vspace{5pt}

\begin{enumerate}
	\item $\Omega$ ist eine Menge mit endlicher Anzahl an Elementen. $|\Omega|= n $
	\item $\mathcal{F}$ ist eine Menge der Ereignisse $\textsce $ und eine Teilmenge von $\Omega$. $\mathcal{F} \subseteq \Omega$
	\item $\mathds{P}$ bezeichnet die Sicherheit der Wahrscheinlichkeit, das Wahrscheinlichkeitsmaß. Dieses weist jedem Ereignis in $\mathcal{F}$ eine reelle Zahl zu.  $\mathds{P}: \mathcal{F} \longrightarrow \mathds{R}$
	
	
\end{enumerate}

\vspace{10pt}



\subsection *{Kolmogorovs Axiome lauten:}

\vspace{5pt}

\begin{enumerate}
	\item Die Wahrscheinlichkeit eines Ereignisses ist eine positive, reelle Zahl. \\ $\mathds{P} (\mathds{E}) \geq 0$ f\"ur alle $\mathds{E} \in \mathcal{F}$.
	\item Die Menge aller Ergebnisse bezeichnet man als sicheres Ereignis, das die Wahrscheinlichkeit $1$ hat. $\mathds{P} (\Omega) = 1$
	\item Die Wahrscheinlichkeit der Vereinigung von disjunkten Ereignisssen ist gleich der Summe der Wahrscheinlichkeiten der disjunkten Ereignisse. 
	\vspace{3pt}
	\begin{equation*}
	\mathds{P} (\cup_{i=1}^m \mathds{E}_{i}) = \sum_{i=1}^n \mathds{P} (\mathds{E}_{i})$ f\"ur alle $\mathds{E}_{1},...,\mathds{E}_{m} \in \mathcal{F}
	\end{equation*}
\end{enumerate}
\vspace{5pt}

$\mathds{E}_{1},...,\mathds{E}_{m}$ ist disjunkt, falls $\mathds{E}_{i} \cap \mathds{E}_{j} = \O$ f\"ur alle $i,j= 1,2,...,m$. 
Hierbei bezeichnet $\O$ ein unm\"ogliches Ereignis $\O \in \mathcal{F}$. \\ Im Gegensatz dazu gibt es das sichere Ereignis mit ${\Omega} \in \mathcal{F}$.

\vspace{5pt}

Die Wahrscheinlichkeit wird errechnet durch $ \mathds{P} = \frac {|\mathds{E}|}{|\Omega|}$.
Alternativ definiert man: 

\vspace{5pt}


$p_{i} \geq 0$  f\"ur $i \in \Omega$, sodass $\sum_{i\in \Omega} p_{i} =1$, $p_{i} = \mathds{P} ({i})$



\vspace{10pt}

\subsection *{Konsequenzen der Axiome:}

\vspace{5pt}

\begin{enumerate}
	\item $\mathds{P} (\O) = 0$
	\item Monotonie: $A\subseteq B$ daraus folgt: $\mathds{P} (A) \leq \mathds{P} (B)$
	\item $0 \leq \mathds{P} (E) \leq 1$
	\item $ \mathds{P} (A \cup B) = \mathds{P} (A) + \mathds{P} (B) - \mathds{P} (A \cap B)$
\end{enumerate}


\vspace{10pt}

\subsection *{Zufallsvariablen}

\vspace{5pt}

In einem Wahrscheinlichkeitsraum $(\Omega, \mathcal{F}, \mathds{P})$ beschreibt die Funktion $X:\Omega \longrightarrow \mathds{R}^k$ einen k-dimensionalen Zufallsvektoren. Im Fall von $k=1$ redet man von einer Zufallsvariablen.

\vspace{10pt}

\subsection *{Erwartungswert}

\vspace{5pt}

 Der Erwartungswert einer Zufallsvariablen beschreibt die Zahl, die die Zufallsvariable im Mittel annimmt.

\vspace{3pt}

\begin{equation*}
\mathds{E} [X] = \sum_{x \in \Omega} X (i) * \mathds{P} ({i}) = \sum_{x\in\mathcal{X}} X* \mathds{P} (X=x) wobei $\mathcal{X} $
\end{equation*}
  ein Wertebereich von x ist  

\vspace{10pt}

\subsection *{Varianz}

\vspace{5pt}

Die Varianz kennzeichnet die Ausdehnung einer Wahrscheinlichkeit. Die Varianz einer Zufallsvariable X.

\vspace{3pt}

\begin{equation*}
X= \mathds{E} [X^2] - \mathds{E} [X]^2 = \mathds{E} [ (X - \mathds{E} [X] )^2 ] 
\end{equation*}

\vspace{10pt}

\subsection *{Unabh\"angigkeit von Ereignissen}

\vspace{5pt}

\begin{enumerate}
	\item Wenn f\"ur Ereignisse $E, F \in \mathcal{F} $ gilt
	\vspace{3pt}
	$\mathds{P} (E\cap F)= \mathds{P} (E) * \mathds{P} (F)$
	\vspace{3pt}
dann sind E und F unabh\"angig. 

\item Wenn f\"ur die Zufallsvariablen 
\begin{equation*}
X_{1}: \Omega_{1} \longrightarrow \mathds{R}, \mathcal{F}_{1} , X_{2}: \Omega_{2} \longrightarrow \mathds{R}, \mathcal{F}_{2}
\end{equation*}
 gilt
\vspace{3pt}
\begin{equation*}
\mathds{P} ({X_{1} \in A} \cap {X_{2} \in B} ) = \mathds{P} ({X_{1} \in A} * \mathds{P} {X_{2} \in B} )
\end{equation*}
\vspace{3pt}
f\"ur alle $A, B \subseteq \mathds{R}$
\end{enumerate}

\vspace{10pt}

\subsection *{Marginal bestimmte Wahrscheinlichkeiten}

\vspace{5pt}

\begin{equation*}
\mathds{P} ({X_{1} \in A } ) = \sum_{b \in \Omega_{2} } \mathds{P} {X_{1} \in A} \cap {X_{2} \in b})
\end{equation*}

\vspace{10pt}

\subsection *{Bedingte Wahrscheinlichkeit}

\vspace{5pt}

\begin{equation*}
\mathds{P} ({X_{1} \in A } | { X_{2} \in B } ) = \frac {\mathds{P} ( {X_{1} \in A } \cap {X_{2} \in B })} {\mathds{P} ({X_{2} \in B})}
\end{equation*}

\vspace{10pt}

\subsection *{Summe von Zufallsvariablen}

\vspace{5pt}
\begin{equation*}
X_{1}:\Omega \longrightarrow \mathds{R}, X_{2}:\Omega \longrightarrow \mathds{R}
\end{equation*}
\begin{equation*}
Y=X_{1}+X_{2}
\end{equation*}
\vspace{3pt}
Annahme: $X_{1}$ und $X_{2}$ sind unabh\"angig.\\
\begin{equation*}
{Y=l}=\cup_{i=-\infty}^{\infty} {X_{1}=1} \cap {X_{2}=l-i}
\end{equation*}
\vspace{3pt}
\begin{equation*}
\mathds{P} ({Y=l}) = \sum_{i=-\infty}^{\infty} \mathds{P} ({X_{1}=i} \cap {X_{2}=l-i}) = \sum_{i=-\intfy}^{\intfy} \mathds{P} ({X_{1} = i}) * \mathds{P} ({X_{2}=l-i})
\end{equation*}

\vspace{10pt}

\subsection *{Theorem: Liniarit\"at des Erwartungswerts}

\vspace{5pt}

Angenommen, wir h\"atten eine Funktion
\vspace{3pt}

$f: \mathds{R} \longrightarrow \mathds{R}$ , f\"ur die gilt: 
\begin{equation*}
\int_{\mathds{R}} f(x) dx = 1
\end{equation*}

\vspace{3pt}

dann k\"onnten wir eine Wahrscheinlichkeitsverteilung folgenderma"sen definieren:
\begin{equation*}
\mathds{P} ({X\in A}) = \int_{X \in A} f(x) dx
\end{equation*}




\section{Optimierung}

\subsection{Grundlagen der Optimierung}
Bei der Optimierung betrachtet man vor allem Funktionen, welche Elemente eines $n$-dimensionalen Vektorraumes auf reelle Zahlen abbilden:

\begin{equation*}
f: \mathbb{R}^n \rightarrow \mathbb{R}
\end{equation*}

Beispielhafte Anwendungszwecke wären die Minimierung oder Maximierung der Kosten bzw. der Einnahmen eines Unternehmens. Diese würden dem reellen Funktionswert der Funktion entsprechen und ergäben sich aus den Eigenschaften des Unternehmens wie Gehalt, Anzahl der Mitarbeiter, Marketing-Ausgaben usw., welche man zusammenfassend als Vektor eines $n$-dimensionalen Vektorraumes darstellen könnte.

\subsection{Regression}

\begin{Thm}[Methode der kleinsten Quadrate]
Beispiel für die Lösung eines solchen Minimierungsproblems wäre die Methode der kleinsten Quadrate. Gegeben ist hierbei eine lineare Funktion $f$ und Punkte $P_i=(x_i,y_i)$, welche näherungsweise auf der Geraden liegen:

\begin{align*}
f(x) &= ax \\
y_i &= f(x_i) + \epsilon_{i}
\end{align*}

\begin{dsafigure}
\begin{center}
\includegraphics[width=0.5\textwidth]{\media Grafik-Optimierung_LineareRegression.pdf}
\caption{Konvexe Funktion}
\label{figure:Grafik-Optimierung_LineareRegression.pdf}
\end{center}
\end{dsafigure}

Wenn nun der Abstand $\epsilon_i$ minimiert werden soll, dann gilt:

\begin{equation*}
\min h(x)=\frac{\sum_{i=1}^n(y_i-ax_i)^2)}{e(y_i,x_i)}
\end{equation*}

Hierbei könnte $e^2,\mid e\mid, h(e)$ exemplarische Möglichkeiten seien. Zur Berechnung dieses Minimums gibt es nun verschiedene Möglichkeiten:

\begin{enumerate}
\item Gewöhnliche Tiefpunktberechnung \\
z.~B.$\qquad h'(a)=0\qquad \wedge\qquad h''(a)>0$
\item Gradientenabstiegsmethode \\
$a^{t+1}  =a^{t} - \lambda \cdot \delta f(a^{t})$ \\
Man findet also ein $\lambda > 0$, sodass
$f(a^{t+1}) < f(a^{t})  \qquad \mid \delta f(a^{(t)}) \neq 0$ \\
$g^{+}(\lambda)=f(a^+-\lambda\Delta f(a^t))$ \\
$\lambda^+=argmin(g^+(\lambda))$
\end{enumerate}

\end{Thm}

\begin{Def}[Lokales und globales Minimum]
Für $f: D\rightarrow \mathbb{R}$ ist $x\in D$ ein lokales Minimum, wenn mindestens eine Umgebung $N$ existiert, sodass \\$\forall_{y\in N}$ gilt: $f(y)>=f(x)$ \\
$x\in D$ ist dann ein globales Minimum, falls $N=D$.
\end{Def}

\subsection{Konvexität}

\begin{Def}[Konvexe Funktion]
Eine Funktion $f: \mathbb{R}^n\rightarrow\mathbb{R}$ ist dann eine konvexe Funktion, falls $\forall_{x,y\in D}$ gilt: 
$f(\lambda x+ (1 - \lambda)y) \leq \lambda f(x)+(a-\lambda)f(y) \qquad \mid \forall_{y\in[0,1]}$

\begin{dsafigure}
\begin{center}
\includegraphics[width=0.45\textwidth]{\media Grafik-Optimierung_KonvexeFunktion.jpg}
\caption{konvexe Funktion}
\end{center}
\end{dsafigure}

\end{Def}

\subsubsection{Beispiel zur Konvexität}

\begin{align*}
f(u) &=\mid\mid u\mid\mid_2^2\qquad\mid u\in\mathbb{R}\\
f(\lambda u+(1-\lambda)v) &=\mid\mid\lambda u+(1-\lambda)v\mid\mid_2
\leq\mid\mid\lambda u\mid\mid_2+\mid\mid(1-\lambda)v\mid\mid_2
\leq\lambda\mid\mid u\mid\mid_2+(1-\lambda)\mid\mid v\mid\mid_2
\end{align*}
\subsection{Konvexe Mengen}

\paragraph{Einführung}

Im Rahmen der Optimierung von konvexen Funktionen ist es erforderlich, den Begriff der konvexen Menge einzuführen. Um diese Thematik anschaulich darzustellen, verwenden wir zunächst verschiedene geometrische Figuren in Abbildung \ref{figure:Grafik-Optimierung_KonvexeMengen}, von denen einige konvex sowie andere wiederum nicht konvex sind. \\
Ist es eine direkte Verbindungsstrecke zwischen zwei beliebigen zu finden, die selbst ebenfalls in der Menge liegt, so ist die Menge konvex.

\begin{dsafigure}
\begin{center}
\includegraphics[width=0.4\textwidth]{\media Grafik-Optimierung_KonvexeMengen.pdf}
\label{figure:Grafik-Optimierung_KonvexeMengen}
\caption{Beispiele konvexer Mengen}
\end{center}
\end{dsafigure}

\begin{Def}[Konvexe Menge]

Eine Menge $X$ heißt konvex, falls für alle $x, y \in X$ und $\lambda \in \mathbb{R}$ und $\lambda \in [0,1]$ gilt:

\begin{equation}
\lambda x + (1 - \lambda) \in X
\end{equation}

\end{Def}

\subsection{Beispiele konvexer Mengen}

\paragraph{Kugel}

Ein Beispiel für eine konvexe Menge ist die Menge $S$ der Vektoren, die eine Kugel mit dem Radius $r = 1$, deren Mittelpunkt im Ursprung liegt, beschreibt:

\begin{equation*}
S = \{x \in \mathbb{R}^{n} | \lVert x \rVert \le 1\}
\end{equation*}

\paragraph{Quadrant}

Auch die Quadranten des kartesischen Koordinatensystems lassen sich durch eine konvexe Menge erfassen. So gilt z.~B. für den ersten Quadranten:

\begin{equation*}
L = \{ x \in \mathbb{R}^{n} | x_1 \ge 0, x_2 \ge 0, ..., x_n \ge 0\}
\end{equation*}

\paragraph{Eistüte}

Bei einer Eistüte bzw. einem quadratischen Kegel handelt es sich um ein weiteres Beispiel für eine konvexe Menge, da sie folgendermaßen für $x \in \mathbb{R}^{n-1}$ und $t \in \mathbb{R}$ beschrieben werden kann:

\begin{equation*}
Q = \{(x, t) \in \mathbb{R}^{n} | \lVert x\rVert ^{2} \le t\}
\end{equation*}

\paragraph{Box}

Die Menge von Vektoren in $B \subseteq \mathbb{R}$ bilden eine konvexe Menge in Form einer Box mit der Seitenlänge 2, deren Mittelpunkt im Ursprung liegt:

\begin{equation*}
B = \{x \in \mathbb{R}^{n} | |x_{i}| \le 1, i = 1, ..., n\}
\end{equation*}

\subsection{Konvexe Programme}

\paragraph{Einführung}

Bei einem konvexen Programm handelt es sich um die Optimierung einer Funktion $f$, die vom Parameter $x$ abhängig ist. Wir versuchen diese Funktion unter Berücksichtigung von Nebenbedingungen zu minimieren. Hierbei sind sowohl die Zielfunktion als auch die Menge der Punkte, die die Nebenbedingungen erfüllt, konvex.

\begin{Def}[Konvexes Programm]
Sei $f: \mathbb{R}^{n} \rightarrow \mathbb{R}$ eine konvexe Funktion und $X \subseteq \mathbb{R}^{n}$, so ist das konvexe Optimierungsproblem bzw. Programm die Minimierung der Funktion $f(x)$ unter der Nebenbedingung, sodass $x \in X$.
\end{Def}

\paragraph{Lineare Programme}

Lineare Programme lassen sich durch lineare Funktionen beschreiben, z.~B. $f(x) = \langle c, x \rangle$. Darüber hinaus sind eine lineare Abbildungsmatrix $A$, z.~B. $A = \begin{pmatrix}1 & 0 \\ 0 & 1 \end{pmatrix}$, sowie ein Vektor $b$ und ein weiterer Vektor $c$, z.~B. $c = \begin{pmatrix}1 \\ 0 \end{pmatrix}$, gegeben. $x$ ist dann zulässig, wenn $x \in X$ mit $X: Ax \le b$ bzw. $X = \{x \in \mathbb{R}^{n} | Ax \le b\}$ erfüllt ist. \\
Ein Beispiel für ein lineares Programm wäre:

\begin{equation*}
f(x) = \langle \begin{pmatrix}1 \\ 0 \end{pmatrix}, \begin{pmatrix}x_{1} \\ x_{2} \end{pmatrix} \rangle = x_{1}
\end{equation*}

Für die Nebenbedingungen gilt:

\begin{equation*}
\begin{pmatrix}1 & 0 \\ 0 & 1 \end{pmatrix}
\end{equation*}

Für den Wertebereich des linearen Programms folgt: $x_{1} \le 1$, $x_{2} \le 1$.

\paragraph{Quadratische Programme}

Im Vergleich zu einem linearen Programm weist z.~B. eine Kostenfunktion im quadratischen Programm eine quadratische Form auf. \\
Beispielsweise sei ein quadratisches Programm $f(x) = x^{T}Qx$ mit $Q = \begin{pmatrix}1 & 0 \\ 0 & 1 \end{pmatrix}$. $f$ gilt es nun unter der Bedingung $x \in X$ für $X: Ax \le b$ zu minmieren:

\begin{align*}
f(x) &= x^{T}Qx \\
&= \begin{pmatrix}x_{1}, x_{2}\end{pmatrix} \begin{pmatrix}1 & 0 \\ 0 & 1\end{pmatrix} \begin{pmatrix}x_{1} \\ x_{2}\end{pmatrix} \\
&= \begin{pmatrix}x_{1}, x_{2}\end{pmatrix} \begin{pmatrix}x_{1} \\ x_{2}\end{pmatrix} \\
&= x_{1}^{2} + x_{2}^{2}
\end{align*}
\subsection{Gradientenmethode}

Folgender Algorithmus implementiert die Gradientenmethode, die einen wichtigen Teil des Gebietes der Optimierung darstellt. Funktionieren tut sie wie folgt:
\begin{algorithmic}[1]

\Procedure{$\mathbf{calculateMinimum}$}{}
   \\$F(x) = (x^2-2)^2$
\\F'(x) = 4x$^3$-8x

\\
\\x = 10.0
\\$\lambda = 0.001$

\For{$i = 1, \dots, m$}
 \State $\lambda = \lambda+0.001\cdot i$
  \For{$j = 1, \dots, n$}
    \State $x = x - \lambda \cdot F'(x)$
  \EndFor
\EndFor
 
 
\\Print x\EndProcedure
\Statex
\end{algorithmic}

\paragraph{Erklärung}
$F(x)$ beschreibt die Funktion, deren globales Minimum wir finden wollen. $F'(x)$ ist dementsprechend die Ableitung der Funktion $F(x)$. Mit $x=10.0$ setzen wir den Schätzwert, ab dem optimiert wird. $\lambda$ bekommt einen niedrigen Wert deklariert, damit man sich in kleinen Schritten dem Minimum annähern kann. Dies wiederum geschieht in zwei For-Schleifen, deren Inhalt in diesem Fall je 8 mal durchlaufen wird. In der äußeren Schleife sorgen wir dafür, dass unser Lambda größere Werte annimmt. In der inneren For-Schleife wird die Schätz-Variable $x$ mit Hilfe der Ableitung und $\lambda$ angepasst. Daraus folgt:
\begin{equation*}
 f(x+\lambda \nabla f(x)) \leq f(x)
 \end{equation*}



%\newcommand{\media}{../../media/}


\section{Machine Learning}
\author{Farhadiba Mohammed, Dennis Kempf, David Steinmann}
Machine Learning ist ein Konzept, bei dem der Computer durch Algorythmen in der Lage ist Probleme selbstständig zu lösen. Dabei erkennt der Computer nach einer gewissen Lernphase Gesetzmässigkeiten, durch die er dann in der Lage ist, neue, ihm unbekannte Datensätze bezüglich eines bestimmten Problems zu lösen.

Dass der Computer die Gesetzmässigkeiten erkennen kann, werden verschiedene Systeme genutzt, doch zunächst einmal sollten wir uns mit den Begriffen Klassifikation und Regression vertraut machen.


\subsection{Klassifikation und Regression}
\author{David Steinmann}
Bei Klassifikation und Regression handelt es sich um zwei verscheidene Arten von Problemen, beziehungsweise Möglichkeiten zur Problemlösung im Bereich Machine Learning.

\subsubsection{Klassifikation}
Bei der Klassifikation sorgt der Algorythmus am Ende dafür, dass der Datensatz klassifiziert wird. Das heißt, der Datensatz wird vom Computer in verschiedene Klassen eingeteilt. Die Klassen müssen vorher vom Mensch entschieden werden. Der Computer klassifiziert die Datensätze dann anhand der Attribute des entsprechenden Datensatzes. Die klassifizierten Datensätze enthalten dann meistens eine weitere Dimension, in der die Klassifizierung anhand einer Zahl gespeichert ist.

\subsubsection{Regression}
Bei der Regression geht es darum, den entsprechenden Datensätzen am Ende einen festen Wert zuzuordnen. Der Unterschied zur Klassifikation besteht darin, dass der zugeordnete Wert nicht für eine Klasse steht, sondern ein konkretes Ergebnis beinhaltet. Ein Beispiel wäre der Preis von einer Wohnung. Diese wurde mit vielen verscheidenen Attributen in die Datenbank eingespeichert (z. B. Größe, Lage ...) der Endwert wäre dann zum Beispiel der konkrete Preis und nicht eine Klassifizierung in "`teuer"', "`mittelteuer"', "`billig"'. 
%Noch die Mathematische Regression (methode der kleinsten Quadrate)

\subsection{Support Vector Machines}
\author{David Steinmann}
Support Vector Machines oder kurz SVMs sind ein Konzept zum Lösen von Machine Learning Problemen.
Dabei werden für alle Datensätze zürst gewisse Features festgelegt. Diese Features entsprechen den Attributen des Datensatzes. Je nach Problem kann es unterschiedlich viele Features geben, doch allgemein kann man sagen, dass, je mehr Features es sind die Genauigkeit des Programms genaür wird.

Diese Features werden normalerweise in einem Vektor für den entsprechenden Datensatz gespeichert.
Dieser Vektor wird $x_{i}$ genannt wobei $i = 1, ..., n$ den Datensatz beschreibt mit  $x \in  \mathbb{R}^m$.
Zusätzlich hat jeder Datensatz noch einen weiteren Wert, welcher das >>Ergebnis<< des Datensatzes beschreibt, der $y_{i} \in \mathbb{R}$ oder auch >>Label<< genannt wird.

Die Features und das Label zusammen beschreiben einen Punkt im n-dimensionalen Koordinatensystem, wobei $n = m + 1$, da zu den m Dimensionen von x noch die eine Dimension von y dazukommt.

Durch die Features ist es dann möglich den Datensätzen einen bestimmten Ort im n-dimensionalen Koordinatensystem zuzuordnen. Diese können dann auf Grund ihrer räumlichen Anordnung klassifiziert werden. Dabei soll der Abstand von beiden Gruppen von Punkten zu dem Trennobjekt maximal sein. Dadurch kann die Klassifizierung optimal durchgeführt werden. Andernfalls würden schon geringe Abweichungen von den Test-Datensätzen reichen, das ein Datensatz falsch klassifiziert würde. (\ref{SVM2}) 


\begin{dsafigure}
\begin{center}
	\includegraphics[width=0.5\textwidth]{\media svm.jpg}
	\caption{Da der Abstand zwischen den unterschiedlich klassifizierten Datensätzen maximiert werden soll, gilt die rote und nicht die blaue Linie als Trennelement.}
	\label{SVM1}
	\end{center}
\end{dsafigure}


Um die Trennlinie zu optimieren, gibt es auch eine mathematische Beschreibung. Da sich Ebenen auch in der Form $\langle x, w \rangle = b $ beschreiben lassen, folgt, dass wir die Datensätze auch etwas anders betrachten können (\ref{SVM2}).
Das daraus folgende mathematische Programm sieht wie folgt aus: 

\begin{align*}
	\operatorname{maximiere}_{w,b} \quad & \frac{2}{\norm{(w)}} \\
	\operatorname{sodass} \quad &y_{i}(\langle w,x_{i} \rangle -b) \geq 1 \quad \forall  i = 1,...,n ; \\
	&x \in \mathbb{R}^n; \quad	y \in \{1;-1\}
	 %\label{SVMProblemMax}
\end{align*}

Unter der Vorraussetzung, dass alle vorhandenen Datensätze richtig klassifiziert sind. Dieses Problem kann 
man aber genauso schreiben als:

\begin{align*}
		\operatorname{minimiere}_{w,b} \quad & \norm{(w)} \\
		\operatorname{sodass} \quad &y_{i}(\langle w,x_{i} \rangle -b) \geq 1 \quad \forall i = 1,...,n ; \\
		&x \in \mathbb{R}^n; \quad y \in \{1; -1\}
		%\label{SVMProblemMin}
\end{align*}
%Genauso gemacht wie obendrüber, funktioniert aber nicht geaunso ??

Klassifikation:

\begin{align*}
	y_{i} =& \{ 1 \operatorname{falls} \langle w,x_{i} \rangle - b \geq 1\} \\
		& \{ -1 \operatorname{falls} \langle w,x_{i} \rangle \leq 1\}
	%\label{KlassifikationSVM}
\end{align*}
%????

Nachdem das Problem optimiert wurde, ist der Computer in der Lage, weitere Daten einzuordnen, vorrausgesetzt, es ist richtig optimiert.
Dennoch kann es bei dieser Art von SVMs zu einigen Problemen kommen.\\

\begin{itemize}
	\item Wenn die Datensätze nicht linear separierbar sind, d.h., es ist nicht möglich die beiden Datensätze mit einem linearen Element zu trennen\\
	
	\item Wenn die Daten nicht alle korrekt klassifiziert sind oder es Daten gibt, die in dem >>Bereich<< der anderen Seite liegen.
\end{itemize}

\subsubsection{Soft Margin SVMs}
Die Soft Margin SVMs können mit dem zweiten Problem umgehen.
Sie sind im Prinzip wie herkömmliche SVMs aufgebaut, nur dass sie eine gewisse Fehlertoleranz besitzen. Diese kommt durch eine Veränderung an der Formel zustande:

\begin{align*}
	\operatorname{minimiere}_{w,b} \quad & \norm{w} + C\sum_{i}{z_{i}} \quad \forall i = 1,...,n \\
	\operatorname{sodass} \quad & y_{i}(\langle w,x_{i} \rangle - b) \geq 1 \quad \forall i = 1,...,n; \\ 
	&x \in \mathbb{R}^n; y \in \{1; -1\} \quad z_{i} \geq 0 \\
	%\label{eq: SoftMarginSVMs}
\end{align*}

Dabei steht $z_{i}$ für die Grösse des Fehlers des Punktes $x_{i}$ und C ist eine Konstante, deren grösse bestimmt, wie stark die Fehler gewichtet werden, da die Konstante mit der Summe der Fehler multipliziert wird. Da die gesamte Gleichung aber minimiert werden soll sorgt ein hohes C dafür, dass das Problem schwerer optimiert werden kann. Die Konstante muss vom Mensch selbst gewählt werden, je nachdem, wie schwer Fehler gewichtet werden sollen.

\begin{dsafigure}
	\begin{center}
		\includegraphics[width=0.5\textwidth]{\media Figure_SVM.pdf}
		\caption{Links und rechts zur Trenngeraden befinden sich die parallelen Grenzen (gestrichelte Geraden). Ziel der Optimierung ist es, den Abstand zwischen den Grenzen zu maximieren, um den Normalenvektor $w$ zu bestimmen.}
		\label{FigSVM}
	\end{center}
\end{dsafigure}

\section{Neuronale Netze}
\author {Farhadiba Mohammed}

\subsection{Einführung}

Inspiriert von unserem Verständnis, wie das menschliche Gehirn lernt, benutzen Neuronale Netze Lernalgorithmen, welche besonders für praktische Anwendungen geeignet sind.
Dazu zählen Spracherkennung, Objekterkennung in Bildern und die Fähigkeit individuell passende Produkte vorzuschlagen, die dem Kunden gefallen könnten. 
Ein Neuronales Netz wird von mehreren Schichten aufgebaut. Ausgangsschicht ist dabei, die Datenschicht, auf die ein oder mehrere hidden layer folgen. Als Ausgabewert erhält man schließlich einen Vektor, welcher die Wahrscheinlichkeitsverteilung darstellt. Mit anderen Worten, wie wahrscheinlich es ist, dass das Ausgangsobjekt zu einer bestimmten Klasse gehört.

\begin{dsafigure}
	\begin{center}
		\includegraphics[width=0.5\textwidth]{\media Figure_NN.pdf}
		\caption{Ein vollständig verbundenes neuronales Netzwerk mit $i$ Eingängen und $k$ Ausgängen, bestehend aus $n$ Schichten mit jeweils $m$ >>Neuronen<<.}
		\label{FigNN}
	\end{center}
\end{dsafigure}

\subsection{Deep Learning}

Mit Deep Learning beschreibt man die Neuronalen Netze, die über mehr als einen >>versteckte Schicht<< verfügen. Damit ist gemeint, dass sich zwischen Ein- und Ausgabeschicht weitere Schichten befinden. 

Um die Anzahl an Parametern zu vergrößern, werden diese auch über nicht-lineare Funktionen miteinander verknüpft. Dadurch kann eine höhere Genauigkeit erzielt werden. Jedoch kann eine zu hohe Anzahl an Parametern auch dazu führen, dass das Netzwerk >>overfitted<<, d.h., zu sehr an den Trainingsdatensatz angepasst, wird. Dann kann das Netzwerk neue, fremde Daten nicht mehr korrekt klassifizieren. 

Wie wir in der Abbildung 2.3 sehen können ist ein typisches Fully connected Neuronales Netz abgebildet. Als Basis findet man unten die Datenschicht mit ihren Eingabewerten in Form eines Vektors $x_1 ... x_i$. Diese Werte werden nun mit einem jeweiligen Faktor $w$ multipliziert. Hierbei handelt es sich um ein Skalarprodukt. Das Ergebnis wird nun im ersten hidden layer abgespeichert. Darauf wird eine nicht lineare Funktion angewendet und das Ganze wird erneut mit einem Faktor $w$ multipliziert. Dieser Vorgang wiederholt sich so lange bis uns schließlich ein Vektor mit seinen Komponenten $p_1 .. p_1$ zurückgegeben wird, welcher als Wahrscheinlichkeitsverteilung interpretiert wird. Anhand eines konkreten Beispiels würde es folgendes bedeuten:
Nehmen wir an wir haben ein Bild und wollen ermitteln zu welcher Klasse Haus, Tisch oder Stuhl das daurauf abgebildete Objekt gehört. Unsere inputs wären demnach die Pixel des Bildes. Diese durchlaufen nun das Neuronale Netz und wir erhalten eine Wahrscheinlichkeitsverteilung als Rückgabewert, welche uns mitteilt, dass es am wahrscheinlichsten ist, dass das abgebildete Objekt ein Objekt der Klasse Stuhl ist. Demnach ist die Klassifizierung vollzogen.


\begin{align*}
a_j^{(k)} = f^{(k)} (\langle w_{ji}^{(k)}, a_i^{(k-1)}\rangle) = f^{(k)} (\sum_{i=1}^{m^{(k-1)}} w_{ji}^{(k)},a_i^{(k-1)} )
\end{align*}

\subsection{Convolutional Neural Networks}

\begin{dsafigure}
	\begin{center}
		\includegraphics[width=0.5\textwidth]{\media Figure_ConvNN}
		\caption{Ein Convolutional Neural Network (CNN) (dt.: >>faltendes neurales Netzwerk<<) mit vier Eingängen ($x_1, ..., x_4$) und zwei Ausgängen ($p_1, p_2$). Dazwischen befindet sich eine Schicht aus drei >>Neuronen<<, die als Filter wirkt.}
		\label{FigConvNN}
	\end{center}
\end{dsafigure}


Bei den Convolutional Networks handelt es sich um ein Neuronales Netz mit einer vereinfachten Struktur. Diese kommt dadurch zu Stande, dass nun kein Fully Connected Neuronales Netzwerk mehr vorliegt. Das bedeutet nicht jeder der Knoten ist mit jedem verbunden, sondern nur jeweils die lokalen drei.


\subsection{SVMs als Neuronale Netze}

Wir haben bereits SVMs kennengelernt. Im folgenden wird verdeutlicht, wie man diese auch als Neuronales Netz ausdrücken kann. Charakteristisch für dieses Neuronale Netz ist, dass es nur einen Layer besitzt:

\begin{align*}
\min\limits_{b,w,z} \norm{w} + C  \sum_{i=1}^{n} z
\end {align*}

sodass:

\begin{align*}
&y_i (<w,x> - b) \geq 1 - z_i ; z_i \geq 0 \\
\implies &z_i \geq 1- y_i (\langle w,x \rangle - b) \text{ und } z_i \geq 0
\end {align*}


\begin{align*}
\min &\norm{w} + C \cdot \sum_{i=1}^{n} \max (0, 1-y_i (\langle w,x_i\rangle-b) \\
 = &\norm{w} + C \sum_{i=1}^{n} \ell ( 1-y_i (\langle w,x_i\rangle)-b) \\
=&\norm{w} + C \sum_{i=1}^{n} \ell (a_i)
\end{align*}


\subsection{Lineare Regression}

Die lineare Regression kann auch in Form von Neuronalen Netzen ausgedrückt werden. 


\begin{align*}
\min\limits_{w,b} 0,5 \sum_{i=1}^{n} ((\langle w,x_i\rangle- b) -y_i)^2
a_i = \langle w,x_i\rangle -b-y_i
\end{align*}




%\nocite{*} % alle referenzen anzeigen, sogar wenn sie nicht im text zitiert sind
%\bibliography{lit}{}
%\bibliographystyle{plain}
\end{document}
